At amplitude 0.8, standard IDEC lost 14.49 percentage points of accuracy relative to clean images. Stripe NMI rose from 0.0004 to 0.1799. Switching orientation changed its predicted cluster on 41.72\% of test images. The projection reduced this to 39.50\% (paired reduction $2.21\pm0.92$ points) and stripe NMI to 0.1658. Both nuisance measures decreased in all five seeds. Mean digit ACC increased by $0.64\pm0.73$ points, with seed differences of $+0.59,+0.76,+0.28,-0.20,+1.76$ points. Figure 3 shows why the mean alone overstates consistency.

The shared initialization already flipped on 42.79\% of strong-stripe inputs. Standard fine-tuning reduced that rate slightly. Removing reconstruction entirely yielded 60.54\% ACC, the highest mean of the three methods, and 39.51\% flips, close to the projection. Projected reconstruction retained higher digit NMI (0.581 versus 0.560). The ablation supports a contribution from reconstruction during fine-tuning, but most of the sensitivity was already present before this stage. At this strength, fine-tuning took roughly 13.4, 14.4 and 9.3 seconds per seed on the RTX 5080 for standard, projected and clustering-only arms, respectively.
